% Prepared for Edward Baker, 23 October 2026.
% Substantial LLM assistance: gpt-6-astra (Codex), reasoning effort ultra.
% Working manuscript contribution; auxiliary details are in Notes 10 and 11.
% Expository clarification, 23 October 2026: gpt-6.1-sol (Codex), effort xhigh.

\section{Collision spectra, positive-time transfer, and the missing global estimate}
\label{sec:spectral-transfer}

For distinct real points $x_1,\ldots,x_N$, let $e_i$ denote the $i$th
standard basis column vector in $\R^N$, and write
\[
 L(x)=\sum_{i<j}\frac{(e_i-e_j)(e_i-e_j)^T}{(x_i-x_j)^2},
 \qquad q(t)=\Tr\!\left(4tL(x(t))[I+4tL(x(t))]^{-1}\right).
\]

\paragraph{Geometric meaning and choice of the filter.}
Let $V(x)=(x_i^{j-1})_{i,j=1}^N$ be the Vandermonde matrix and put
$\Delta=(\det V)^2=\prod_{i<j}(x_i-x_j)^2$. Direct differentiation gives
\begin{equation}
 L=-\operatorname{Hess}_x\log|\det V|
   =-\tfrac12\operatorname{Hess}_x\log\Delta,
 \qquad
 v^TLv=\sum_{i<j}\frac{(v_i-v_j)^2}{(x_i-x_j)^2}.
 \label{eq:vandermonde-hessian}
\end{equation}
Thus $L$ is the curvature matrix of the logarithmic repulsion energy,
and also the weighted graph Laplacian with inverse-square gap weights.
The Hessian here is the full matrix of second derivatives with respect to
the root positions; the scalar spatial Laplacian would give its trace.
The constant vector is a zero mode because a common translation changes
none of the gaps.

For $t>0$ on a real simple slice, let $\lambda_k(t)\ge0$ be the eigenvalues
of $L(x(t))$. Then
\begin{equation}
 q(t)=\sum_{k=1}^N f(4t\lambda_k(t)),
 \qquad f(u)=\frac{u}{1+u}.
 \label{eq:filter-effective-dimension}
\end{equation}
This is a smooth count, or effective dimension, of the modes at the heat
scale: a mode with $4t\lambda_k\ll1$ contributes little, a mode with
$4t\lambda_k\gg1$ contributes nearly one, and a mode at $4t\lambda_k=1$
contributes one half. Each mode contributes at most one, so a single
exceptionally small gap cannot give an arbitrarily large contribution
through one eigenvalue.

The factor $4t$ is calibrated by the elementary heat birth
$H(t,z)=z^2-2t$. Its roots are $\pm\sqrt{2t}$, their squared gap is $8t$,
and their two-point Laplacian has nonzero eigenvalue $2/(8t)=1/(4t)$.
This double therefore contributes exactly $1/2$ to $q$.
The proposition below gives the corresponding spectrum for every local
multiplicity. For a fixed finite population of distinct real endpoint
roots, by contrast, $L(x(t))$ stays bounded and $q(t)\to0$ as $t\downarrow0$.
At positive time, very close descendants of distinct endpoint roots can
also contribute; $q$ is therefore used to bound multiplicity.

The particular function $f$ is a convenient choice of smooth filter.
Besides boundedness, it has the matrix order property
$f(A)=I-(I+A)^{-1}$ for $A\succeq0$, used in the comparison proof below,
and a determinant representation:
\begin{equation}
 q(t)=\left.
 a\,\partial_a\log\det\!\left(I+aL(x(t))\right)
 \right|_{a=4t}.
 \label{eq:filter-log-determinant}
\end{equation}
Here the root configuration and time are held fixed in $\partial_a$;
this is not the total derivative along the heat flow.
Thus the filter is also the logarithmic scale derivative of the determinant
of the regularized curvature matrix. This identity leads to the
forest-polynomial representation in Section~\ref{sec:forest-filter}.

A selected population is \emph{complete at its endpoint sites} if it includes
all multiplicity labels at every site it selects. Its finite matrix contains
only those labels, but their motion is always that of the complete field.
For endpoint multiplicities $m_a$, distinguish
\[
 M_0=\sum_{m_a\ge2}m_a,\qquad D_0=\sum_a(m_a-1),\qquad N=\sum_a m_a.
\]
Thus $M_0$ counts labels at nonsimple sites; $D_0$ counts excess multiplicity.

\begin{proposition}[Endpoint spectrum and scalar multiplicity encoding]
\label{prop:endpoint-spectrum}
Let a finite complete endpoint population of a real-type analytic solution of
$H_t=-H_{zz}$ have real sites, with real simple descendants for small $t>0$.
Then
\[
 \operatorname{spec}(4tL(x(t)))\longrightarrow
 \bigsqcup_a\{0,1,\ldots,m_a-1\},
 \qquad
 q_0:=\lim_{t\downarrow0}q(t)=\sum_a(m_a-\mathsf H_{m_a}),
 \quad \mathsf H_m=\sum_{j=1}^m\frac1j.
\]
In particular $D_0/2\le q_0\le D_0$ and $M_0\le4q_0$.
\end{proposition}
\begin{proof}
The local heat expansion gives
$x_{a,j}(t)=a+\sqrt{2t}\,h_{m_a,j}+O(t)$, where the $h_{m,j}$ are the
zeros of the probabilists' monic Hermite polynomial $\mathrm{He}_m$.
The Hermite equation implies
$\sum_{j\ne i}(h_i-h_j)^{-1}=h_i/2$.
Acting with $L(h)$ on evaluations of the monomial $x^k$ gives
$(k/2)h_i^k$ plus lower-degree terms. The evaluations of degrees
$0,\ldots,m-1$ form a Vandermonde basis, so
$\operatorname{spec}L(h)=\{0,1/2,\ldots,(m-1)/2\}$.
Inter-site entries, including their diagonal contributions, vanish after
multiplication by $t$. The block limits therefore give the stated spectrum.
Applying $u\mapsto u/(1+u)$ yields
$\sum_{k=1}^{m-1}k/(k+1)=m-\mathsf H_m$.
Each positive summand lies between $1/2$ and $1$; also
$m-\mathsf H_m\ge m/4$ for $m\ge2$, with equality at $m=2$.
\end{proof}

The Hermite spectrum is inherited from the local analysis in Note 06; its
scalar trace is an exact encoding, not a theta upper estimate. A double
contributes $1/2$, a triple $7/6$, and a quadruple $23/12$.
These are fixed-population limits, with no uniform exchange of height and
time limits asserted.

\begin{theorem}[Sharp positive-time transfer]
\label{thm:sharp-transfer}
Suppose, in addition, that all adjacent descendants within each selected
endpoint cluster remain globally adjacent in a complete real simple system
on $(0,T]$, and that their gaps satisfy
\[
 (g^2)'\le8,\qquad g(0)=0.
\]
Then $M_0\le4q(t)$ for every $0<t\le T$. The constant four is sharp.
\end{theorem}
\begin{proof}
Every such gap obeys $g(t)^2\le8t$. Partition each even cluster into adjacent
pairs, and each odd cluster of size at least three into pairs and one
consecutive triple. Retain only the pair edges and the two consecutive edges
in each triple. Then $4tL\succeq A_*$, where the pair and triple blocks of
$A_*$ are one half of the respective path Laplacians.
The map $f(A)=I-(I+A)^{-1}$ is operator monotone on positive semidefinite
matrices. A pair block has spectrum $\{0,1\}$ and contributes $1/2$ to
$\Tr f(A_*)$. A triple has spectrum $\{0,1/2,3/2\}$ and contributes
$1/3+3/5=14/15\ge3/4$. Thus each block pays at least one quarter of its
labels. Singleton sites need no payment. For sharpness,
$H(t,z)=z^2-2t$ gives $M_0=2$ and $q(t)=1/2$.
\end{proof}

For canonical theta, the hypothesis on gaps is supplied by the full
canonical real-root motion. For globally adjacent roots $a<b$, $g=b-a$,
\[
 (g^2)'=8-4g^2\sum_{r\ne a,b}\frac1{(a-r)(b-r)}\le8.
\]
The exterior sum is part of the hypothesis and is absolutely convergent;
its terms are nonnegative by global adjacency. No external factor is set to
one. In the conditional theta application, RH supplies the required complete
real system at positive time. The finite matrix need not contain the exterior
roots, but a truncated motion law is insufficient. The improvement from the
earlier constant six concerns proportional transfer; the sufficient
all-simple threshold remains $q(t)<1/2$.

\begin{theorem}[Bounded spectral mass and short adjacent gaps]
\label{thm:gap-filter}
For arbitrary $y_1<\cdots<y_N$, set $A=L(y)$ and
\[
 B_N(\delta)=\frac1N\#\{i<N:y_{i+1}-y_i<\delta\},\qquad
 Q_N(s)=\frac1N\Tr\!\left(sA(I+sA)^{-1}\right).
\]
For $s,\delta>0$,
\begin{equation}
 B_N(\delta)\frac{s}{\delta^2+2s}
 \le Q_N(s)\le2B_N(\delta)+\frac{\pi^2s}{3\delta^2}.
 \label{eq:gap-filter-two-sided}
\end{equation}
Consequently, along any sequence of configurations and positive scales,
\[
 Q_N(s_N)\longrightarrow0
 \quad\Longleftrightarrow\quad
 B_N(a\sqrt{s_N})\longrightarrow0\quad\hbox{for every fixed }a>0.
\]
\end{theorem}
\begin{proof}
Mark every vertex touching a short adjacent gap. If $r$ vertices are marked,
then $r\le2NB_N(\delta)$. The unmarked vertices are $\delta$-separated;
every marked point is also at distance at least $\delta$ from each
unmarked point. Otherwise the unmarked point's nearest original neighbor
on that side would be less than $\delta$ away.

Let $D$ be the principal submatrix of $A$ on the unmarked set. Its diagonal
\emph{retains} interactions with marked vertices. For an unmarked vertex,
its $k$th unmarked neighbor on either side is at distance at least $k\delta$.
For a marked vertex, its $k$th unmarked neighbor on either side satisfies the
same bound. Summing $2\sum_{k\ge1}(k\delta)^{-2}$ first over the unmarked
vertices and then over marked-to-unmarked interactions gives
\[
 \Tr D\le\frac{\pi^2N}{3\delta^2}.
\]
Interlacing pays at most $r$ for the discarded eigenvalues; the remaining
filter values are bounded by $s$ times the eigenvalues of $D$. This proves
the upper bound, including the empty-principal-matrix case.

The short adjacent edges form disjoint paths. From $K$ such edges select a
matching of at least $K/2$ edges. Each retained edge has a Laplacian
eigenvalue at least $2/\delta^2$. Dropping all other edges and using operator
monotonicity proves the lower bound. Substituting $\delta=a\sqrt{s}$ gives
\[
 B_N(a\sqrt{s})\le(a^2+2)Q_N(s),\qquad
 Q_N(s)\le2B_N(a\sqrt{s})+\frac{\pi^2}{3a^2}.
\]
The equivalence follows by taking a limsup at fixed $a$, then letting
$a\to\infty$ for the reverse direction.
\end{proof}

For example, independently acquired bounds
$B_N(\delta)\le C\delta^\beta+\varepsilon_N$, $\beta>0$, imply
\[
 Q_N(s)\le\left(2C+\frac{\pi^2}{3}\right)s^{\beta/(\beta+2)}
             +2\varepsilon_N
\]
whenever $s^{1/(\beta+2)}$ lies in their certified range. This is a sufficient
source input, not a bound proved here for theta.

To state the global interface precisely, let an endpoint target have $N_h$
labels, of which $n_h=N_h-r_h$ are covered by complete clusters. At a
source-prescribed time $t_h$, set $y_i=\nu_hx_i(t_h)$ and
$s_h=4t_h\nu_h^2$. Under Theorem~\ref{thm:sharp-transfer},
\begin{equation}
 \frac{M_h}{N_h}\le\frac{r_h}{N_h}
       +4\frac{n_h}{N_h}Q_{n_h}(s_h).
 \label{eq:global-spectral-interface}
\end{equation}
Thus density-one simplicity would follow conditionally from $r_h=o(N_h)$
and a vanishing fraction of gaps smaller than $2a\sqrt{t_h}$ for every fixed
$a>0$. That microscopic gap estimate remains unpaid. Choosing times after
resolving the unknown endpoint gaps would merely reformulate the target.
The published late-clock onset $x\ge e^{C/t}$ in \cite{polymath} does not
supply this estimate. With its natural unfolding
$\nu_X=g_x(X,t)\sim\log X/(4\pi)$, the uniformly bounded normalized gaps
instead imply $Q_N(s)\ge1-1/N-(1+4s/C_g^2)^{-1/2}$ by path comparison.
For this unfolding, $s=4t\nu_X^2\to\infty$ in the high-height onset regime,
so the unmodified filter saturates. This does not exclude
matched backward transport on a different clock.

\section{A coefficient-level filter with explicit conditioning}
\label{sec:forest-filter}

The preceding filter is defined through roots, but its acquisition need not
isolate them individually. Let $P(z)=\prod_{i=1}^N(z-x_i)$ be monic and
$\Delta=\prod_{i<j}(x_i-x_j)^2$. The following identity is algebraic; real
rootedness is needed only for the positivity conclusions.

\begin{proposition}[Forest polynomial and coefficient enclosure]
\label{prop:forest-filter}
For $a>0$ and distinct roots define $D=\Delta\det(I+aL)$ and $G=a\partial_aD$.
Then
\begin{equation}
 D=\sum_{F\,\operatorname{forest}}
     a^{|E(F)|}\prod_{C\in\operatorname{comp}(F)}|C|
     \prod_{\{i,j\}\notin E(F)}(x_i-x_j)^2,
 \qquad q_a=\frac GD.
 \label{eq:forest-polynomial}
\end{equation}
The sum runs over all spanning labeled forests, including isolated vertices.
The quotient identity is asserted where $D\ne0$. Both $D$ and $G$ are
symmetric polynomials in the roots and hence polynomials in the coefficients
of $P$. Fix $q_*=G_*/D_*$ and form the coupled polynomial $K=G-q_*D$.
If coefficient data give $|D-D_*|\le E_D<|D_*|$ and $|K|\le E_K$, then
\begin{equation}
 |q_a-q_*|\le\frac{E_K}{|D_*|-E_D}.
 \label{eq:forest-residual-bound}
\end{equation}
The separate allowance $E_K=\epsilon_G+|q_*|\epsilon_D$, obtained from
$|G-G_*|\le\epsilon_G$, is a valid fallback.
\end{proposition}
\begin{proof}
Apply Cauchy--Binet to the weighted incidence matrix in $L$.
A cyclic edge selection has dependent columns and contributes zero.
For a forest, its incidence Gram determinant is $\prod_C|C|$: its nonzero
maximal minors omit one vertex per component and have square one.
Thus its determinant contribution is the edge-weight product times
$a^{|E|}\prod_C|C|$. Each edge occurs at most once, so multiplying
by the single product $\Delta$ cancels every denominator. This proves the
polynomial identity. Logarithmic differentiation of the determinant gives
\[
 a\partial_a\log D=a\Tr\bigl(L(I+aL)^{-1}\bigr)=q_a.
\]
The discriminant is independent of $a$. Finally $q_a-q_*=K/D$, and
$|D|\ge|D_*|-E_D$, proving \eqref{eq:forest-residual-bound}.
\end{proof}

The coefficient residuals in the proposition are finite, computable quantities.
For example, on a coefficient polydisc $|c_j-c_{*,j}|\le r_j$, take
$D_*=D(c_*)$, $G_*=G(c_*)$ and use the finite Taylor majorants
$E_D=\sum_{\alpha\ne0}|\partial^\alpha D(c_*)|r^\alpha/\alpha!$
and
$E_K=\sum_{\alpha\ne0}|\partial^\alpha K(c_*)|r^\alpha/\alpha!$,
since $K(c_*)=0$. Form $K$ before taking absolute values, so correlated
coefficient changes can cancel. The estimating gain is that the paid
margin is for $D$, not necessarily for $\Delta$ or individual roots. It is
not a statement that all coefficient balls are well conditioned.

\begin{lemma}[Real repeated-root boundary]
For fixed $a>0$ and a real root multiset, the polynomial extension has $D>0$
if and only if every root multiplicity is at most two.
\end{lemma}
\begin{proof}
All forest terms are nonnegative. A nonzero term must include every edge
joining equal roots, since any omitted such edge contributes zero. A cluster
of three or more forces a cycle and makes every term vanish. If all clusters
have size at most two, the disjoint collapsed pair edges form a forest whose
complementary gaps are nonzero, so its term is positive.
\end{proof}

The distinction is visible already in degrees two and three. For
$P=z^2+bz+c$, put $\delta=b^2-4c$. Then
\[
 D=\delta+2a,\qquad G=2a,\qquad q_a=\frac{2a}{\delta+2a}.
\]
In particular, $P=z^2+p$, $-a/8\le p\le0$, gives $4/5\le q_a\le1$,
including the repeated reference polynomial. There is no discriminant floor
or internal root isolation in this enclosure. It is a stable filter
calculation, not a simplicity certificate: its values are above the relevant
threshold. For complex $|p|\le a/8$, one similarly has
$|D|\ge3a/2$ and $|q_a-1|\le1/3$.

For $P=z^3+bz^2+cz+d$, let $B=b^2-3c$. Directly from the three-vertex forest
sum,
\[
 D=\Delta+2aB^2+6a^2B,\qquad G=2aB^2+12a^2B.
\]
On the real-rooted cone, $0\le\Delta\le4B^3/27$. Thus, for $B>0$ and
$\kappa=B/a$,
\[
 \frac{2\kappa+12}{4\kappa^2/27+2\kappa+6}
 \le q_a\le\frac{2\kappa+12}{2\kappa+6}.
\]
The restricted real limit at a triple is two. There is no corresponding
universal extension in complex coefficient space: along $P=z^3-r^3$,
$B=0$, $\Delta=-27r^6$, and $q_a=0$, whereas the real Hermite approach gives
$q_a\to2$ at fixed $a$.

The heat scale must also be retained. For the cubic heat polynomial
$P_t=z^3-6tz$,
\[
 D(t,a)=108t(a+2t)(a+4t),\qquad
 D(t,4t)=5184t^3,\qquad G(t,4t)=6048t^3,
\]
so $q(t)=7/6$. More generally, the coefficient snapshots
\[
 P_t=z^3-6tz+\eta t^{3/2},\qquad
 q(t)=\frac{6048}{5184-27\eta^2}.
\]
For $\eta\ne0$ these snapshots do not form a heat solution.
For real $\eta$ with $|\eta|\le\sqrt{32}$ the polynomials are real rooted
and $7/6\le q\le7/5$. A constant coefficient perturbation of order $t^{3/2}$
therefore has an order-one effect on the limiting filter. To recover the
Hermite value, sufficient depressed-coefficient accuracies are
$p+6t=o(t)$ and $d_{\rm depressed}=o(t^{3/2})$; the cost of centering must
be included. Likewise the quadratic heat polynomial retains the scale
$\epsilon/t$ for its constant coefficient error. Exact coefficient transport
removes rootwise inverse-gap factors, but cannot remove these source precision
requirements. Notes 10 and 11 give further controls and the associated
complete finite replays.
